\documentclass[12pt,a4paper]{article}

\usepackage[T2A]{fontenc}
\usepackage[utf8]{inputenc}
\usepackage[russian]{babel}
\usepackage{amsmath,amssymb}
\usepackage{booktabs}
\usepackage{geometry}
\usepackage{array}
\usepackage{hyperref}
\usepackage{enumitem}

\geometry{margin=2.5cm}
\hypersetup{
    colorlinks=true,
    linkcolor=black,
    urlcolor=blue
}

\title{Отчёт о проделанной работе\\
по исследованию и реализации предобуславливателей}
\author{}
\date{\today}

\begin{document}

\maketitle
\tableofcontents
\newpage

\section{Введение}

Целью работы было исследовать влияние различных предобуславливателей на скорость и устойчивость решения линейных систем
\[
Ax=b,
\]
а также подготовить основу для построения обучаемого, в перспективе нейросетевого, предобуславливателя.

Работа развивалась в несколько этапов. Сначала был проведён теоретический анализ литературы по итерационным методам, спектральным свойствам матриц и классическим схемам предобуславливания. Затем были реализованы и проверены несколько базовых предобуславливателей, построен единый контур вычислительных экспериментов, выполнен сравнительный анализ по числу итераций, обусловленности и времени. После этого был сделан переход к data-driven подходу: вручную, на \texttt{NumPy}, был реализован обучаемый аппроксиматор обратного оператора. Итогом стала оформленная исследовательская библиотека \texttt{preconditioner} с центральным интерфейсом \texttt{workbench.py}.

\paragraph{Замечание о реконструкции истории.}
Доступная git-история репозитория фиксирует уже в основном упакованный результат: релиз \texttt{v0.1}, демонстрационный пример и ноутбук с онбордингом. Поэтому ранний этап ручных экспериментов на \texttt{NumPy} и литературный обзор в данном отчёте описываются как содержательная реконструкция по текущему коду и по постановке задачи, а не как покоммитная хроника разработки.

\section{Постановка задачи}

При решении больших или плохо обусловленных систем прямое обращение матрицы либо слишком дорого, либо неустойчиво. Поэтому используется предобусловленная система
\[
M^{-1}Ax=M^{-1}b,
\]
где матрица $M$ выбирается так, чтобы:
\begin{enumerate}[label=\arabic*.]
    \item оператор $M^{-1}A$ имел более благоприятный спектр и меньшую обусловленность;
    \item применение $M^{-1}$ к вектору было существенно дешевле, чем решение исходной задачи без предобуславливания.
\end{enumerate}

В рамках работы решались следующие подзадачи:
\begin{enumerate}[label=\arabic*.]
    \item изучить классические подходы к предобуславливанию;
    \item реализовать набор предобуславливателей и итерационных решателей;
    \item организовать серию воспроизводимых экспериментов на разных семействах матриц;
    \item ввести количественные метрики сравнения;
    \item перейти от фиксированных аналитических схем к обучаемому предобуславливателю;
    \item свести все компоненты в единый исследовательский фреймворк.
\end{enumerate}

\section{Анализ литературы и мотивация}

На этапе литературного обзора были важны следующие линии рассуждений.

\subsection{Итерационные методы и спектральная картина}

Для методов семейства Krylov, таких как GMRES, LGMRES и BiCGSTAB, характер сходимости тесно связан со спектром матрицы и её обусловленностью. Если собственные значения разнесены слишком широко, либо матрица близка к вырожденной, число итераций растёт, а устойчивость вычислений падает.

\subsection{Классические предобуславливатели}

Классический обзор предобуславливания приводит к идее, что хороший предобуславливатель должен быть компромиссом между качеством приближения матрицы $A$ и стоимостью применения. Отсюда естественно возникают диагональные, факторизационные, структурные и приближённо-обратные схемы.

\subsection{Переход к обучаемым схемам}

Если для некоторого класса матриц удаётся накопить выборку $\{A_i\}$, то можно ставить задачу аппроксимации отображения
\[
A \mapsto A^{-1}
\quad \text{или} \quad
A \mapsto T(A),
\]
где $T(A)$ является удобным приближением обратного оператора. Такая постановка естественным образом связывает предобуславливание с методами машинного обучения. В логике данной работы именно это и было целевым направлением: от аналитических предобуславливателей перейти к нейросетевому или, шире, обучаемому предобуславливателю.

\section{Реализованные математические объекты}

\subsection{Семейства тестовых матриц}

Для корректного сравнительного анализа в проекте были реализованы несколько семейств матриц:
\begin{itemize}[leftmargin=*]
    \item случайные плотные матрицы;
    \item строго диагонально доминирующие матрицы;
    \item симметричные положительно определённые матрицы вида
    \[
    A=B^\top B + \alpha I;
    \]
    \item матрицы Гильберта
    \[
    H_{ij}=\frac{1}{i+j-1},
    \]
    используемые как заведомо плохо обусловленный тест;
    \item матрицы Тёплица;
    \item матрицы со случайным ортогональным базисом и контролируемым спектром
    \[
    A=Q\Lambda Q^\top,
    \]
    что позволяет целенаправленно задавать уровень обусловленности.
\end{itemize}

Такой набор покрывает как неструктурированные, так и структурированные случаи, а также даёт возможность тестировать методы на заведомо тяжёлых примерах.

\subsection{Метрики качества}

В проекте использовались как стандартные, так и специальные метрики.

\paragraph{Число обусловленности.}
Для общей матрицы рассматривалось стандартное число обусловленности
\[
\kappa_2(A)=\|A\|_2\|A^{-1}\|_2.
\]

Для симметризованной версии в коде используется signed-метрика:
\[
A_s=\frac{A+A^\top}{2}, \qquad
\kappa_s(A)=\frac{\lambda_{\max}(A_s)}{\lambda_{\min}(A_s)}.
\]
Такая величина удобна в исследовательском режиме, когда важно отслеживать не только масштаб спектра, но и возможную смену знака минимального собственного значения.

\paragraph{Изменение обусловленности после предобуславливания.}
После построения предобуславливателя вычисляется матрица
\[
A_{\text{pre}} = M^{-1}A,
\]
и затем сравниваются:
\[
\Delta_\kappa = \kappa(A_{\text{pre}})-\kappa(A),
\qquad
\rho_\kappa = \frac{\kappa(A_{\text{pre}})}{\kappa(A)}.
\]

\paragraph{Итерационные и временные метрики.}
Помимо спектральных характеристик, измеряются:
\begin{itemize}[leftmargin=*]
    \item число итераций до остановки метода;
    \item норма невязки
    \[
    \|r_k\|_2 = \|b-Ax_k\|_2;
    \]
    \item время построения предобуславливателя;
    \item время оценки обусловленности;
    \item время решения системы;
    \item суммарное время полного пайплайна.
\end{itemize}

\section{Реализация классических предобуславливателей}

В библиотеке реализован базовый набор предобуславливателей, достаточный для сравнительного анализа и для построения сильных baseline-моделей.

\subsection{Отсутствие предобуславливания}

Базовая модель \texttt{None} используется как контрольная точка, относительно которой удобно оценивать эффект всех остальных подходов.

\subsection{Диагональный предобуславливатель}

Если
\[
A = D + R,
\]
где $D$ состоит из диагональных элементов матрицы, то простейшая схема Якоби использует
\[
M=D,
\qquad
M^{-1} = D^{-1}.
\]
Этот подход очень дёшев, легко реализуется и часто заметно улучшает поведение для диагонально доминирующих или близких к ним систем.

\subsection{LU-предобуславливатель}

Если доступна факторизация
\[
A \approx LU,
\]
то в качестве предобуславливателя можно применять решение двух треугольных систем:
\[
M^{-1}x = U^{-1}L^{-1}x.
\]
В плотном случае такой предобуславливатель может быть очень сильным и по сути близок к точному решению, но его построение дорого, поэтому в коде он ограничен по размеру матрицы.

\subsection{Циркулянтный предобуславливатель}

Для структурированных матриц полезна циркулянтная аппроксимация. В проекте она строится через усреднение диагоналей, после чего применение $M^{-1}$ выполняется через дискретное преобразование Фурье:
\[
C = \operatorname{circ}(c), \qquad
C = F^\ast \Lambda F, \qquad
C^{-1} = F^\ast \Lambda^{-1} F.
\]
Это делает метод особенно интересным для матриц с выраженной почти-тёплицевой структурой.

\subsection{SVD-предобуславливатель}

Для матрицы
\[
A = U\Sigma V^\top
\]
можно построить псевдообратный оператор
\[
A^+ = V\Sigma^+ U^\top.
\]
Такой предобуславливатель близок к приближённому обратному, но вычислительно дорог, поэтому в исследовании он также используется как эталонный, а не массовый вариант.

\section{Итерационные методы решения}

Для оценки предобуславливателей реализованы и подключены следующие итерационные схемы:
\begin{itemize}[leftmargin=*]
    \item GMRES;
    \item LGMRES;
    \item BiCGSTAB.
\end{itemize}

Все методы оформлены единообразно: фиксируются число итераций, траектория невязок, факт сходимости и итоговая ошибка решения. Это позволяет сравнивать не только качество самого предобуславливателя, но и его совместимость с конкретным типом решателя.

\section{Экспериментальный контур и сравнительный анализ}

\subsection{Серийные эксперименты}

В проекте выделен отдельный слой для постановки массовых вычислительных экспериментов. Эксперименты выполняются по сетке размеров матриц, по набору семейств и по списку предобуславливателей и солверов. Для каждого запуска формируются записи с полным набором метрик, а затем строятся агрегированные таблицы средних значений и стандартных отклонений.

Это позволило перейти от единичных ручных тестов к систематическому сравнению методов. По сути, именно на этом этапе работа приобрела исследовательский, а не просто прикладной характер.

\subsection{Сравнение по размеру и по обусловленности}

В коде реализованы два основных сценария:
\begin{enumerate}[label=\arabic*.]
    \item sweep по размеру матрицы $n$;
    \item sweep по заранее заданным уровням $\kappa(A)$ при фиксированном размере.
\end{enumerate}

Второй сценарий особенно важен, поскольку позволяет исследовать не только абсолютный масштаб задачи, но и то, как методы ведут себя при ухудшении спектральной картины.

\subsection{Визуализация результатов}

Для анализа поведения методов предусмотрено построение графиков:
\begin{itemize}[leftmargin=*]
    \item зависимость изменения обусловленности от размера матрицы;
    \item зависимость числа итераций от размера матрицы;
    \item зависимость метрик обусловленности от исходного $\kappa(A)$;
    \item зависимость времени решения и полного времени пайплайна от размера;
    \item зависимость времени от уровня обусловленности.
\end{itemize}

Таким образом, итог исследования выражается не только в численных таблицах, но и в удобной визуальной форме, пригодной для интерпретации.

\section{Обучаемый предобуславливатель как шаг к нейросетевому подходу}

\subsection{Мотивация}

Содержательная цель работы состояла не только в сравнении классических схем, но и в попытке построить предобуславливатель, который \emph{обучается} на примерах матриц и затем автоматически выдаёт приближение к обратному оператору.

В текущем состоянии репозитория этот этап представлен модулем \texttt{ml.py}. Формально он ещё не реализует глубокую нейронную сеть, однако уже воплощает саму идею обучаемого предобуславливателя и был написан вручную на \texttt{NumPy}, без внешних высокоуровневых ML-фреймворков. Это важный исследовательский этап: сначала была проверена сама постановка задачи и стабильность пайплайна обучения, а уже затем можно переходить к более сложной архитектуре.

\subsection{Постановка задачи обучения}

Для обучающей выборки строятся пары
\[
(A_i, T_i),
\]
где в роли целевого оператора могут выступать:
\begin{itemize}[leftmargin=*]
    \item точная обратная матрица $A_i^{-1}$;
    \item псевдообратная матрица $A_i^+$;
    \item диагональное приближение обратной матрицы.
\end{itemize}

Матрица $A_i$ разворачивается в вектор признаков
\[
\phi(A_i)=\operatorname{vec}(A_i),
\]
при необходимости с добавлением bias-компоненты.

Обучаемая модель ищет оператор $W$, минимизирующий функционал гребневой регрессии:
\[
\min_W \sum_{i=1}^{N}\|W\phi(A_i)-\operatorname{vec}(T_i)\|_2^2 + \lambda \|W\|_F^2.
\]

После обучения предсказание строится как
\[
\widehat{T}(A)=\operatorname{mat}(W\phi(A)),
\]
и далее используется в качестве предобуславливателя:
\[
M^{-1}(A)\approx \widehat{T}(A).
\]

\subsection{Почему это связано с нейросетевым предобуславливателем}

Хотя реализованная модель линейна, она уже соответствует общей идее параметризованного оператора
\[
P_\theta(A),
\]
где параметры $\theta$ извлекаются в ходе обучения. С этой точки зрения текущий модуль можно интерпретировать как baseline для будущего нейросетевого предобуславливателя:
\begin{itemize}[leftmargin=*]
    \item постановка задачи обучения уже сформулирована;
    \item генерация датасета уже реализована;
    \item интерфейс превращения модели в предобуславливатель уже есть;
    \item сравнение с классическими методами уже встроено в общий фреймворк.
\end{itemize}

То есть переход к настоящей нейронной сети потребует скорее замены внутренней модели, чем переделки всей исследовательской инфраструктуры.

\section{Формирование исследовательского фреймворка \texttt{workbench.py}}

\subsection{Причина появления фреймворка}

По мере роста числа матриц, предобуславливателей, солверов и графиков ручной сценарий экспериментов становится неудобным. Возникает типичная исследовательская проблема: вычислительный код уже есть, но он размазан по функциям и одноразовым скриптам, а значит трудно расширяется и плохо воспроизводится.

Именно поэтому в проекте сформировался центральный модуль \texttt{workbench.py}. Его задача состоит в том, чтобы собрать разрозненные эксперименты в один воспроизводимый API.

\subsection{Что делает \texttt{PreconditionerWorkbench}}

Класс \texttt{PreconditionerWorkbench} предоставляет:
\begin{itemize}[leftmargin=*]
    \item дефолтные семейства матриц;
    \item дефолтные предобуславливатели;
    \item дефолтные итерационные методы;
    \item запуск sweep-экспериментов;
    \item оценку одной матрицы;
    \item стандартный набор графиков для анализа результатов.
\end{itemize}

Тем самым он закрывает типичный исследовательский цикл: сгенерировать матрицу, построить предобуславливатель, решить систему, измерить метрики, агрегировать статистику и визуализировать результат.

\subsection{Что добавляет \texttt{PreconditionerWorkbenchV2}}

Вторая версия фреймворка решает уже инженерную задачу расширяемости. В ней реализованы реестры компонентов:
\begin{itemize}[leftmargin=*]
    \item регистрация новых семейств матриц;
    \item регистрация новых предобуславливателей;
    \item регистрация новых солверов;
    \item возможность передавать компонент по имени, объектом или функцией-конструктором.
\end{itemize}

Это очень важный результат работы, потому что после него проект перестаёт быть набором демонстрационных примеров и становится платформой для дальнейших исследований.

\subsection{Архитектура итогового пакета}

В оформленном виде проект распадается на следующие модули:

\begin{center}
\begin{tabular}{>{\raggedright\arraybackslash}p{4cm} >{\raggedright\arraybackslash}p{9cm}}
\toprule
\textbf{Модуль} & \textbf{Назначение} \\
\midrule
\texttt{matrices.py} & генерация тестовых семейств матриц и сеток размеров \\
\texttt{preconditioners.py} & реализация аналитических предобуславливателей \\
\texttt{solvers.py} & единый интерфейс к GMRES, LGMRES и BiCGSTAB \\
\texttt{metrics.py} & вычисление числа обусловленности и производных метрик \\
\texttt{experiments.py} & проведение серий экспериментов и агрегация результатов \\
\texttt{plotting.py} & построение графиков для сравнительного анализа \\
\texttt{ml.py} & обучаемый аппроксиматор обратного оператора \\
\texttt{workbench.py} & высокоуровневая точка входа в исследовательский пайплайн \\
\texttt{onboard.ipynb} & демонстрационный ноутбук, показывающий полный цикл использования \\
\bottomrule
\end{tabular}
\end{center}

\section{Наблюдаемые результаты по текущему репозиторию}

По доступному ноутбуку и тестовым прогонам можно зафиксировать несколько конкретных результатов.

\subsection{Демонстрационный sweep}

В ноутбуке \texttt{onboard.ipynb} выполнен быстрый экспериментальный режим со следующими параметрами:
\begin{itemize}[leftmargin=*]
    \item размеры матриц: $n \in \{16, 44, 72, 100, 128\}$;
    \item число повторов на точку: $2$;
    \item число исследуемых семейств матриц: $3$;
    \item предобуславливатели: \texttt{None}, \texttt{Diagonal}, \texttt{LU};
    \item решатель: \texttt{GMRES}.
\end{itemize}

В этом режиме было сформировано $90$ детальных записей экспериментов и $45$ агрегированных записей.

\subsection{Пример одиночной оценки}

Для одной SPD-матрицы размера $32 \times 32$ диагональный предобуславливатель дал следующий демонстрационный результат:
\begin{itemize}[leftmargin=*]
    \item исходная signed-обусловленность: примерно $4956.61$;
    \item обусловленность после предобуславливания: примерно $-129.50$;
    \item отношение обусловленностей: примерно $-2.61 \cdot 10^{-2}$;
    \item решение методом GMRES сошлось за $32$ итерации;
    \item норма невязки составила порядка $7.3 \cdot 10^{-14}$.
\end{itemize}

Этот пример не исчерпывает всей картины, но показывает, что даже простое диагональное предобуславливание способно существенно изменить спектральные характеристики задачи.

\subsection{Результат обучаемой модели}

На обучающей выборке из $48$ SPD-матриц размера $16 \times 16$ обучаемый аппроксиматор обратного оператора показал:
\begin{itemize}[leftmargin=*]
    \item среднюю относительную ошибку по норме Фробениуса порядка $8.38 \cdot 10^{-7}$;
    \item среднее число обусловленности предобусловленной системы порядка $1.00023$.
\end{itemize}

Важно подчеркнуть, что это демонстрационный внутренний тест на обучающей выборке. Тем не менее, он подтверждает главную идею: отображение ``матрица $\rightarrow$ приближённый обратный оператор'' действительно можно учить и затем использовать как предобуславливатель.

\section{Основные результаты работы}

По итогам выполненной работы можно сформулировать следующие содержательные результаты:
\begin{enumerate}[label=\arabic*.]
    \item Выполнен теоретический анализ проблемы предобуславливания и выделены ключевые классы методов.
    \item Реализованы несколько классических предобуславливателей: диагональный, LU, циркулянтный и SVD-подход.
    \item Построен единый интерфейс для запуска итерационных методов решения.
    \item Организован воспроизводимый контур массовых экспериментов на различных семействах матриц.
    \item Введены метрики сравнения по обусловленности, невязке, числу итераций и времени.
    \item Реализована обучаемая модель аппроксимации обратного оператора, написанная вручную на \texttt{NumPy}.
    \item Подготовлена инфраструктура, в которой обучаемый предобуславливатель можно сравнивать с классическими методами на общих основаниях.
    \item Сформирован итоговый исследовательский фреймворк \texttt{workbench.py}, который объединяет генерацию данных, построение предобуславливателей, запуск экспериментов и визуализацию.
\end{enumerate}

\section{Заключение и направления развития}

Проведённая работа показывает последовательный переход от теоретического анализа и ручных численных экспериментов к оформленной исследовательской платформе. Сначала были изучены свойства классических предобуславливателей и реализованы их базовые варианты. Затем был организован честный сравнительный анализ на контролируемых семействах матриц. После этого была сделана ключевая попытка перейти от фиксированных аналитических формул к обучаемому предобуславливателю, реализованному вручную на \texttt{NumPy}. Наконец, все наработки были сведены в единый фреймворк, пригодный для дальнейшего расширения.

С практической точки зрения наиболее важным итогом является не только набор уже написанных методов, но и сама архитектура исследования. Теперь проект готов к следующему шагу: замене текущего линейного обучаемого аппроксиматора на полноценную нейросетевую модель, например MLP или другую архитектуру, способную лучше учитывать структуру матрицы и обобщаться на более широкий класс задач.

В качестве естественных продолжений работы можно выделить:
\begin{itemize}[leftmargin=*]
    \item переход от линейной ridge-модели к глубокой нейросети;
    \item исследование обобщающей способности на тестовой, а не только обучающей выборке;
    \item работу с более крупными и разреженными матрицами;
    \item добавление неполных факторизаций и специализированных предобуславливателей;
    \item автоматический экспорт графиков и таблиц в формат полноценного отчёта.
\end{itemize}

Таким образом, проделанная работа уже образует законченный исследовательский цикл и одновременно служит прочной базой для следующего шага, а именно для построения полноценного нейросетевого предобуславливателя.

\end{document}
